\section{Tessellations}

The substrate was chosen by comparison, not assumption, and the comparison rests on a full catalog.

\textbf{The catalog.} Of the regular hyperbolic honeycombs up to five dimensions, forty-five are catalogued and forty-two are buildable by the integer engine, the remaining three being star or apeirogonal. A single reusable measurement battery was run across all forty-two.

\textbf{What is general and what is specific.} The framework ports to every buildable substrate, the reversible charge-conserving rule, exact conservation, the finite light cone, the Bethe-lattice gravity, the holography, the de Sitter cosmology, and the radial-tree hierarchy, and a propagating fermion appears on all forty-two. What is selective is the native spinor structure of the 24-cell, which appears on only seven, the honeycombs whose docks are $[3,4,3]$-faceted, namely $\{3,4,3,4\}$, $\{4,3,4,3\}$, and five in five dimensions.

\textbf{The discriminators.} The three-dimensional $\{5,3,4\}$, with its twelve icosahedral directions, carries no spinor, its twelve directions splitting as $1+3+3'+5$. They are not a root system, so it has no gauge group, and its cusp is two-dimensional, so its gravity is logarithmic rather than $1/r$, which makes it the decisive control. The two-dimensional $\{7,3\}$ is thinner still, with a one-dimensional cusp, and serves as the cleanest stage for holography. Among the seven spinor-carriers, $\{3,4,3,4\}$ alone is crystallographic, hyperbolic in its bulk, and flat and three-dimensional in its cusp at once. Its five-dimensional cousin the pentacomb $\{3,4,3,3,4\}$ is the one that carries both the spinor and the curvature the committed substrate trades away, the richer relative held in reserve.
